\chapter*{Abstract}
\addcontentsline{toc}{chapter}{Abstract}
\begin{SingleSpace}
The exponential growth of digital educational resources has precipitated a paradigm shift in autonomous learning. However, this abundance of unstructured multimedia data—ranging from protracted video lectures to dense academic PDFs—imposes an acute cognitive burden on university students. Addressing the fundamental constraint of information overload, this thesis introduces Lectura: a comprehensive, AI-driven digital study assistant engineered to metamorphose diverse, unstructured content into highly consolidated and interactive learning artifacts. By synthesizing state-of-the-art Natural Language Processing (NLP) with pedagogical frameworks like Bloom's Taxonomy, the Spaced Repetition System (SRS), and Multimedia Learning objectives, Lectura fundamentally accelerates knowledge acquisition and retention.

The platform architecture is instantiated as a resilient, cloud-native full-stack web application. The frontend leverages React 18, TypeScript, and Tailwind CSS to deliver a highly reactive user interface featuring inline WYSIWYG editing workflows. Concurrently, the backend is architected in Go (Golang), utilizing a highly concurrent worker pool to manage long-running AI inference cycles, with data persistence managed via PostgreSQL and Redis. The core generative intelligence is powered by Google’s Gemini 3 Flash Large Language Model. It executes executing complex, templated prompts to yield multi-format summaries, context-aware quizzes, algorithmic flashcards, localized conversational tutoring, and—crucially—visually immersive, Gamma-standard generative presentations natively rendered and exportable to high-fidelity PDF and PPTX formats.

This document elucidates the theoretical grounding, architectural blueprint, and technical implementation of the Lectura ecosystem. It critically evaluates contemporary literature on AI in education (AIEd), automated visual instructional design, and cognitive retrieval strategies. Furthermore, a rigorous comparative analysis of incumbent platforms (Gamma.app, Notion AI, Quizlet) justifies the necessity for a unified, multimodal ingestion pipeline. By delineating the API ontology, data chronologies, dynamic slide rasterization mechanics, and synchronous worker mechanisms, this project empirically demonstrates how advanced generative paradigms can be elegantly codified into robust software engineering practice, ultimately delivering an uncompromising, highly performant academic workstation.
\end{SingleSpace}

\chapter*{Аңдатпа}
\addcontentsline{toc}{chapter}{Аңдатпа}
\begin{SingleSpace}
Цифрлық білім беру ресурстарының экспоненциалды өсуі автономды оқытуда парадигманың өзгеруіне әкелді. Дегенмен, сағаттық бейнелекциялардан бастап күрделі академиялық PDF файлдарына дейінгі құрылымдалмаған мультимедиялық деректердің көптігі университет студенттеріне үлкен когнитивтік жүктеме түсіреді. Ақпараттық шамадан тыс жүктеменің іргелі мәселесін шеше отырып, бұл дипломдық жұмыс Lectura жобасын ұсынады: әртүрлі, құрылымдалмаған мазмұнды жоғары топтастырылған және интерактивті оқу артефактілеріне айналдыру үшін жасалған жасанды интеллектке негізделген кешенді цифрлық оқу көмекшісі. Заманауи Табиғи тілді өңдеу (NLP) әдістерін Блум таксономиясы, Интервалдық қайталау жүйесі (SRS) және Мультимедиялық оқыту сияқты педагогикалық құрылымдармен біріктіре отырып, Lectura білімді меңгеру мен есте сақтауды түбегейлі жеделдетеді.

Платформа архитектурасы сенімді, бұлтқа бағытталған толық стек веб-қосымшасы ретінде іске асырылған. Фронтенд inline WYSIWYG өңдеу жұмыс үрдістері бар пайдаланушы интерфейсін ұсыну үшін React 18, TypeScript және Tailwind CSS технологияларын пайдаланады. Баэкенд Go (Golang) тілінде жасалған және ұзақ ЖИ генерациялау циклдерін басқаруға арналған жоғары параллельді жұмысшы пулын (worker pool) пайдаланады. Негізгі генеративті интеллект Google-дың Gemini 3 Flash Үлкен тілдік моделі (LLM) арқылы жұмыс істейді. Ол көп форматалы конспектілерді, мәнмәтінге негізделген тесттерді, алгоритмдік флэш-карталарды, чат-репетиторлықты және ең бастысы — тікелей экранда көрсетілетін және PDF, PPTX форматтарына экспортталатын «Gamma» стандартындағы визуалды презентацияларды генерациялайды.

Бұл құжат Lectura экожүйесінің теориялық негіздерін, архитектуралық жобасын және техникалық іске асырылуын түсіндіреді. Онда білім берудегі ЖИ (AIEd), визуалды дизайнды автоматтандыру және когнитивті іздеу стратегиялары бойынша қазіргі заманғы әдебиеттер бағаланады. Қолданыстағы платформаларды (Gamma.app, Notion AI, Quizlet) қатаң салыстырмалы талдау біріктірілген мультимодальды конвейердің қажеттілігін негіздейді. Платформа динамикалық слайд растерлеу механикасын және синхронды жұмысшы механизмдерін талдай отырып, алдыңғы қатарлы ЖИ алгоритмдерінің жоғары деңгейлі бағдарламалық инженерияға қалай интеграцияланатынын көрсетеді.
\end{SingleSpace}

\chapter*{Аннотация}
\addcontentsline{toc}{chapter}{Аннотация}
\begin{SingleSpace}
Экспоненциальный рост цифровых образовательных ресурсов привел к смене парадигмы в автономном обучении. Однако изобилие неструктурированных мультимедийных данных — от многочасовых видеолекций до сложных академических PDF-файлов — налагает огромную когнитивную нагрузку на студентов университетов. Решая фундаментальную проблему информационной перегрузки, данная дипломная работа представляет Lectura: комплексного цифрового учебного ассистента на базе ИИ, разработанного для преобразования разнообразного неструктурированного контента в высококонсолидированные и интерактивные учебные артефакты. Путем синтеза современных методов обработки естественного языка (NLP) с педагогическими концепциями, такими как Таксономия Блума, Система интервальных повторений (SRS) и Мультимедийное обучение, Lectura фундаментально ускоряет усвоение и запоминание знаний.

Архитектура платформы реализована как отказоустойчивое, облачное полностековое веб-приложение. Фронтенд использует React 18, TypeScript и Tailwind CSS для обеспечения высокореактивного пользовательского интерфейса, оснащенного инлайн-редактором (WYSIWYG) для работы с контентом прямо на странице. Бэкенд спроектирован на Go (Golang) и использует высокопараллельный пул воркеров (worker pool) для управления длительными циклами логического вывода ИИ. Ядром генеративного интеллекта выступает большая языковая модель Google Gemini 3 Flash. Она выполняет сложные структурированные запросы для создания многоформатных конспектов, контекстно-зависимых тестов, алгоритмических флэш-карточек, диалогового репетиторства и, что критически важно, — визуально захватывающих генеративных презентаций в стандарте «Gamma», которые нативно рендерятся и экспортируются в форматы PDF и PPTX с высокой точностью (high-fidelity).

В этом документе разъясняются теоретические основы, архитектурный проект и техническая реализация экосистемы Lectura. Проводится критическая оценка современной литературы по ИИ в образовании (AIEd), автоматизированному визуальному проектированию (instructional design) и стратегиям когнитивного извлечения памяти. Строгий сравнительный анализ существующих платформ (Gamma.app, Notion AI, Quizlet) доказывает необходимость единого конвейера мультимодального ввода. Описывая механику динамической растеризации слайдов и механизмы синхронных пулов, проект эмпирически демонстрирует, как передовые нейросетевые парадигмы могут быть интегрированы в надежную практику программной инженерии, предоставляя бескомпромиссную академическую рабочую станцию.
\end{SingleSpace}
